% cau-hoi.tex: gói 2 bộ GAME Jelly Trail Academy (jelly). Module 3 String & Date (9 câu), Module 4 Aggregation (9 câu), case c3, c4.
% SQL nằm trong solution.sql, nối với câu bằng mã câu. Biên dịch: XeLaTeX. Hướng dẫn: HUONG_DAN.md.
\documentclass[11pt,a4paper]{article}
\usepackage{aio-sql-cauhoi}
% Vá XeLaTeX (cong-cu/va-xelatex.tex), không sửa .sty.
\setmonofont{NotoSansMono}[Extension=.ttf, UprightFont=*-Regular, BoldFont=*-Bold, Scale=0.86]
\renewcommand{\erdkhoa}[1]{\node[font=\small, color=cqMuted, anchor=north west, align=left, text width=14cm] at ([yshift=-1mm]current bounding box.south west) {\textbf{Khoá:} #1};}

% ================================================================ KHAI BÁO DÙNG CHUNG
% ---------------------------------------------------------------- bộ dữ liệu
\khaibaodulieu{jelly}{GAME · Jelly Trail Academy}
  {Game giải đố ghép kẹo trên di động (iOS, Android, 8 thị trường): người chơi, phiên chơi, lượt chơi màn, mua trong ứng dụng, quảng cáo, chi phí thu hút; 60 ngày từ 2026-07-01 tới 2026-08-29.}

% ---------------------------------------------------------------- mảng nghiệp vụ
\khaibaodomain{customer}{Người chơi}{Hoạt động, giữ chân, phễu hướng dẫn, nhóm cài đặt.}
\khaibaodomain{sales}{Mua trong ứng dụng}{Giao dịch IAP, gói vật phẩm, người chơi trả tiền.}
\khaibaodomain{marketing}{Thu hút và sự kiện}{Kênh, chiến dịch, chi phí, sự kiện LiveOps, thử nghiệm A/B.}
\khaibaodomain{operations}{Vận hành game}{Màn chơi, độ khó, booster, phiên chơi, phiên bản.}
\khaibaodomain{finance}{Doanh thu}{Tổng doanh thu, ARPDAU, ARPPU, CPI, ROAS.}
\khaibaodomain{after_sales}{Chất lượng sau phát hành}{Sự cố, hoàn tiền, đối soát sau bản vá.}

% ---------------------------------------------------------------- khái niệm
% \khaibaokhainiem{ma}{Tên}{Định nghĩa}{Công thức}{Trong bộ dữ liệu}
\khaibaokhainiem{dau}{DAU (số người chơi hoạt động trong ngày)}
  {DAU là số người chơi khác nhau có mở game ít nhất một lần trong một ngày.}
  {số user\_id khác nhau có hoạt động trong ngày}
  {bảng \code{player_daily}: mỗi dòng là một người chơi trong một ngày có hoạt động; \code{daily_kpi.dau} là số tổng hợp sẵn, khớp với cách đếm này.}
\khaibaokhainiem{nhom-cai-dat}{Nhóm cài đặt (cohort)}
  {Nhóm cài đặt là tập người chơi cài game trong cùng một khoảng thời gian, được theo dõi theo số ngày kể từ ngày cài.}
  {}
  {\code{users.install_date}; \code{player_daily.cohort_day} = số ngày từ ngày cài tới ngày hoạt động, D0 là ngày cài.}
\khaibaokhainiem{phien-choi}{Phiên chơi và thời lượng phiên}
  {Phiên chơi là một lần người chơi mở game cho tới khi rời đi. Thời lượng phiên là số phút của phiên đó.}
  {thời lượng phiên bình quân = tổng duration\_minutes / số phiên}
  {\code{sessions}; mã phiên có hai dạng, phiên hướng dẫn có \code{entry_screen} = \code{'tutorial'}.}
\khaibaokhainiem{ty-le-thang-man}{Tỷ lệ thắng màn}
  {Tỷ lệ thắng màn là phần lượt chơi hoàn thành mục tiêu của màn trên tổng số lượt chơi màn đó.}
  {số lượt result = 'complete' / tổng số lượt}
  {\code{level_attempts}; mẫu số gồm cả lượt thua, bỏ cuộc và gặp sự cố.}
\khaibaokhainiem{ty-le-su-co}{Tỷ lệ sự cố trên 1.000 phiên}
  {Tỷ lệ sự cố đo số lần ứng dụng tự thoát do lỗi trên mỗi 1.000 phiên chơi.}
  {tổng crash\_count / tổng sessions × 1000}
  {\code{player_daily.crash_count}, \code{player_daily.sessions}. Số lượt chơi có \code{result} = \code{'crash'} là nguồn đếm khác, không khớp với \code{crash_count}.}
\khaibaokhainiem{phien-ban-luc-choi}{Phiên bản lúc chơi và lúc cài}
  {Phiên bản lúc chơi là phiên bản ứng dụng người chơi dùng trong ngày hoạt động; người chơi có thể cập nhật sau khi cài.}
  {}
  {\code{player_daily.app_version}, \code{sessions.app_version}, \code{level_attempts.app_version} là phiên bản lúc chơi; \code{users.app_version_install} là phiên bản lúc cài. Bản 2.1.0 có từ 30/7, 2.1.1 từ 15/8.}
\khaibaokhainiem{doanh-thu-iap}{Doanh thu IAP}
  {Doanh thu IAP là tiền người chơi trả để mua gói vật phẩm ảo trong ứng dụng.}
  {tổng price\_usd của giao dịch không hoàn tiền}
  {bảng \code{purchases}: mỗi dòng là một giao dịch, giá theo USD; trong kỳ chưa có giao dịch hoàn tiền (\code{is_refund} = 0).}
\khaibaokhainiem{doanh-thu-quang-cao}{Doanh thu quảng cáo}
  {Doanh thu quảng cáo là tiền mạng quảng cáo trả cho game khi người chơi xem quảng cáo.}
  {tổng revenue\_usd}
  {\code{ad_events} (từng lần xem, 5 chữ số thập phân); \code{player_daily.ad_revenue} là tổng theo ngày đã làm tròn 4 chữ số.}
\khaibaokhainiem{nguoi-choi-tra-tien}{Người chơi trả tiền}
  {Người chơi trả tiền là người có ít nhất một giao dịch mua trong ứng dụng trong kỳ xét.}
  {số user\_id khác nhau có giao dịch}
  {\code{purchases}; một người có thể có nhiều giao dịch.}
\khaibaokhainiem{arppu}{ARPPU}
  {ARPPU là doanh thu IAP bình quân trên một người chơi trả tiền.}
  {doanh thu IAP / số người chơi trả tiền}
  {\code{purchases}; khác ARPU, có mẫu số là mọi người chơi.}
\khaibaokhainiem{kenh-thu-hut}{Kênh thu hút và chi phí quảng cáo}
  {Kênh thu hút là nguồn đưa người chơi tới cài game, gồm kênh tự nhiên không trả phí và các kênh quảng cáo trả phí. Chi phí quảng cáo là tiền studio trả cho các kênh trả phí.}
  {}
  {\code{users.channel}, \code{users.campaign_id} ghi kênh lúc cài; kênh tự nhiên \code{'Organic'} có mã \code{'ORG-000'} và không có chi phí. \code{campaign_spend}: mỗi dòng là chi phí một ngày của một chiến dịch tại một quốc gia trên một nền tảng.}

\begin{document}
{\LARGE\bfseries Gói 2 · GAME Jelly Trail Academy · String \& Date, Aggregation\par}
\medskip
{\color{cqMuted}Bộ dữ liệu \texttt{jelly} · Module 3 (9 câu), Module 4 (9 câu), 2 case · chuẩn \texttt{aio-sql-cauhoi} v0.4\par}
\tableofcontents


% ################################################################ MODULE 3 · STRING & DATE

% ================================================================ S1 · write · Cơ bản
\begin{cauhoi}[ma=jelly-s1-gio-cao-diem, dulieu=jelly, domain=operations,
  hinhthuc=write, cap=basic, dang=A3, kynang={HOUR, GROUP BY}, bay={gio-tu-cot-ngay, quen-can-tren},
  tacgia={Agent gói 2}, trangthai=review, bokiemtra={public, hidden-1}, thutu=co]
\tieude{Người chơi mở game nhiều nhất vào giờ nào?}
\boicanh[Anh Bảo · Trưởng nhóm LiveOps]{Anh đang chọn giờ gửi thông báo mời người chơi vào sự kiện mới. Em xem giúp anh trong tháng 7/2026 người chơi hay mở game vào những giờ nào.}
\begin{yeucau}
  \item Xét các phiên chơi bắt đầu từ ngày 1/7/2026 tới hết ngày 31/7/2026.
  \item Đếm số phiên theo giờ bắt đầu phiên (phần giờ của \code{session_start}, từ 0 tới 23).
  \item Trả về \code{gio}, \code{so_phien}; sắp \code{gio} tăng dần.
\end{yeucau}
\begin{erd}
  \erdbang{sessions}{0,0}{session_id:PK, user_id:FK, session_start*=2026-07-05 20:43:00, event_date=2026-07-05, duration_minutes, entry_screen}
\end{erd}
\begin{nghiepvu}{Lên lịch gửi thông báo và mở sự kiện vào khung giờ đông người chơi.}
  \khainiem{phien-choi}
  \luuy{Một người chơi có thể mở game nhiều lần trong ngày, nên số phiên theo giờ cho biết lúc nào game đông, chứ không phải số người chơi khác nhau.}
\end{nghiepvu}
\goiy{Cột ngày và cột thời điểm của phiên khác nhau ở đâu? Xem giá trị mẫu trong ERD.}
\goiy{\code{HOUR(session_start)} trả về phần giờ; lọc tháng 7 có cả cận dưới và cận trên.}
\ghichuthietke{Cơ bản: 1 bảng có giá trị mẫu cho thấy \code{event_date} không có phần giờ. Hai bài sai rớt ở public: lấy giờ từ \code{event_date} (mọi phiên vào giờ 0); thiếu cận trên lẫn phiên tháng 8.}
\end{cauhoi}

% ================================================================ S2 · write · Cơ bản
\begin{cauhoi}[ma=jelly-s2-tach-the-gioi, dulieu=jelly, domain=operations,
  hinhthuc=write, cap=basic, dang=A1, kynang={SUBSTRING_INDEX}, bay={tach-sai-dau-phan-cach},
  tacgia={Agent gói 2}, trangthai=review, bokiemtra={public, hidden-1}, thutu=co]
\tieude{Tên thế giới trong ngữ cảnh màn hình}
\boicanh[Anh Khoa · Thiết kế màn chơi]{Anh đang rà nhãn hiển thị của các màn khó trước bản cập nhật tới. Em lấy giúp anh tên thế giới ghi trong ngữ cảnh màn hình để anh so với tên thế giới gốc.}
\begin{yeucau}
  \item Xét các màn có nhãn độ khó khác \code{'normal'}.
  \item Ngữ cảnh màn hình \code{screen_context} gồm ba phần, ngăn cách bằng chuỗi \code{' > '} (dấu cách, dấu lớn hơn, dấu cách). Lấy phần thứ hai làm \code{the_gioi}, bỏ dấu cách ở hai đầu.
  \item Trả về \code{level_id}, \code{the_gioi}, \code{world_name}; sắp \code{level_id} tăng dần.
\end{yeucau}
\begin{erd}
  \erdbang{level_metadata}{0,0}{level_id:PK, world_name*='Jelly Jungle', screen_context*='Map screen > Jelly Jungle > Level 29', difficulty_tag*='hard', moves_limit}
\end{erd}
\begin{nghiepvu}{Kiểm tra nhãn hiển thị của màn khó trước khi phát hành bản cập nhật.}
  \luuy{Ngữ cảnh màn hình là chuỗi chữ do ứng dụng ghi, gồm màn hình, thế giới và màn. Dấu phân cách có dấu cách ở hai bên, nên tách thiếu dấu cách sẽ để thừa khoảng trắng trong tên.}
\end{nghiepvu}
\goiy{Phần thứ hai nằm giữa dấu phân cách thứ nhất và thứ hai. Cắt hai lần. Lần đầu lấy phần trước dấu phân cách thứ hai, lần sau lấy phần sau dấu phân cách cuối của chuỗi vừa cắt.}
\goiy{\code{SUBSTRING_INDEX(chuỗi, ' > ', 2)} rồi \code{SUBSTRING_INDEX(…, ' > ', -1)}.}
\ghichuthietke{Cơ bản: 1 bảng có giá trị mẫu. Hai bài sai rớt ở public: tách theo \code{'>'} không kèm dấu cách (thừa khoảng trắng hai đầu); chỉ cắt một lần (còn \code{'Map screen > …'}).}
\end{cauhoi}

% ================================================================ S3 · write · Dễ
\begin{cauhoi}[ma=jelly-s3-ngay-trong-tuan, dulieu=jelly, domain=customer,
  hinhthuc=write, cap=easy, dang=A3, kynang={DAYNAME, WEEKDAY, COUNT DISTINCT}, bay={dayofweek-bat-dau-chu-nhat, so-ngay-moi-thu-khac-nhau},
  tacgia={Agent gói 2}, trangthai=review, bokiemtra={public, hidden-1}, thutu=co]
\tieude{Người chơi cài game vào thứ mấy nhiều nhất?}
\boicanh[Chị Mai · Quản lý sản phẩm]{Chị định xếp ca trực hỗ trợ người chơi mới theo ngày trong tuần. Em xem giúp chị từ đầu tháng 7 tới nay, bình quân mỗi thứ trong tuần có bao nhiêu người cài game.}
\begin{yeucau}
  \item Xét người chơi cài từ ngày 1/7/2026 tới hết ngày 29/8/2026.
  \item Tính cho từng thứ trong tuần của ngày cài: số ngày có người cài, số người cài và số người cài bình quân mỗi ngày (số người cài chia số ngày).
  \item Trả về \code{thu} (tên thứ tiếng Anh như MySQL trả về, \code{'Monday'} … \code{'Sunday'}), \code{so_ngay}, \code{so_nguoi_cai}, \code{tb_moi_ngay} (người, làm tròn 1 chữ số).
  \item Sắp theo thứ tự trong tuần, từ thứ Hai (\code{'Monday'}) tới Chủ nhật (\code{'Sunday'}).
\end{yeucau}
\begin{erd}
  \erdbang{users}{0,0}{user_id:PK, install_date*, platform, country, channel}
  \erdbang{player_daily}{1,0}{user_id:FK, activity_date, install_date, cohort_day}
  \erdnoi[kieu=1-n, tren={\code{users.user_id = player_daily.user_id}}, duoi={một người chơi, nhiều ngày hoạt động}]{users.user_id}{player_daily.user_id}
\end{erd}
\begin{nghiepvu}{Xếp ca trực hỗ trợ người chơi mới theo ngày trong tuần.}
  \khainiem{nhom-cai-dat}
  \luuy{Mỗi người chơi có đúng một dòng trong \code{users}, ghi ngày cài. Từ 1/7 tới 29/8/2026 có thứ xuất hiện 9 lần, có thứ chỉ 8 lần, nên so tổng số người cài giữa các thứ sẽ thiên về các thứ có nhiều ngày hơn.}
\end{nghiepvu}
\goiy{\code{DAYOFWEEK} đánh số 1 là Chủ nhật. Hàm nào đánh số 0 là thứ Hai? Mỗi thứ có bao nhiêu ngày trong khoảng xét?}
\goiy{Gom theo \code{DAYNAME(install_date)} và \code{WEEKDAY(install_date)}; số ngày là \code{COUNT(DISTINCT install_date)}.}
\ghichuthietke{Dễ: 2 bảng có điều kiện nối; câu chỉ cần \code{users}. v2 thêm số ngày và bình quân mỗi ngày (theo REVIEW-B): đếm thô thì thứ Bảy cao nhất, bình quân mỗi ngày thì thứ Hai cao nhất. Hai bài sai rớt ở public: sắp theo \code{DAYOFWEEK} (Chủ nhật đứng đầu); chia mọi thứ cho 9.}
\end{cauhoi}

% ================================================================ S4 · write · Dễ
\begin{cauhoi}[ma=jelly-s4-ngay-tu-cai-den-mua, dulieu=jelly, domain=sales,
  hinhthuc=write, cap=easy, dang=A1, kynang={DATEDIFF, JOIN, LIMIT}, bay={dao-thu-tu-datediff, thieu-quy-tac-hoa},
  tacgia={Agent gói 2}, trangthai=review, bokiemtra={public, hidden-1}, thutu=co]
\tieude{Người chơi mua Starter Pack sau khi cài bao lâu?}
\boicanh[Chị Hương · Quản lý doanh thu]{Starter Pack được thiết kế cho người mới, nhưng chị nghe có người chơi lâu rồi vẫn mua. Em lấy giúp chị 20 giao dịch Starter Pack xa ngày cài nhất.}
\begin{yeucau}
  \item Xét các giao dịch có \code{product_type} = \code{'starter_pack'}.
  \item Tính \code{so_ngay} là số ngày từ ngày cài của người chơi (\code{users.install_date}) tới ngày giao dịch (\code{purchases.event_date}), số không âm.
  \item Trả về \code{transaction_id}, \code{user_id}, \code{install_date}, \code{event_date}, \code{so_ngay}.
  \item Sắp \code{so_ngay} giảm dần; bằng nhau thì \code{transaction_id} nhỏ hơn đứng trước. Chỉ lấy 20 dòng đầu.
\end{yeucau}
\begin{erd}
  \erdbang{purchases}{0,0}{transaction_id:PK, user_id:FK, event_date*, product_type*, price_usd}
  \erdbang{users}{1,0}{user_id:PK, install_date*, platform, channel}
  \erdnoi[kieu=n-1, tren={\code{purchases.user_id = users.user_id}}, duoi={nhiều giao dịch, một người chơi}]{purchases.user_id}{users.user_id}
\end{erd}
\begin{nghiepvu}{Quyết định có giới hạn Starter Pack cho người chơi mới trong vài ngày đầu hay không.}
  \khainiem{doanh-thu-iap}
  \luuy{Ngày cài nằm ở bảng người chơi, ngày mua nằm ở bảng giao dịch. Ngày giao dịch luôn từ ngày cài trở đi.}
\end{nghiepvu}
\goiy{\code{DATEDIFF(a, b)} trả về a trừ b. Ngày nào phải đứng trước để số ngày không âm?}
\goiy{Nhiều giao dịch có cùng số ngày ở gần dòng 20. Thứ tự phụ khi bằng số ngày quyết định giao dịch nào nằm trong 20 dòng.}
\ghichuthietke{Dễ: 2 bảng có điều kiện nối. Hai bài sai rớt ở public: đảo đối số \code{DATEDIFF} (ra số âm, lấy nhầm giao dịch gần ngày cài nhất); đảo chiều tiêu chí phụ (nhiều giao dịch bằng số ngày ở ranh giới dòng 20).}
\end{cauhoi}

% ================================================================ S5 · write · Dễ
\begin{cauhoi}[ma=jelly-s5-phien-ban-dang-chuoi, dulieu=jelly, domain=operations,
  hinhthuc=write, cap=easy, dang=A1, kynang={SUBSTRING_INDEX, CAST, GROUP BY}, bay={tach-sai-phan-giua, quen-can-tren},
  tacgia={Agent gói 2}, trangthai=review, bokiemtra={public, hidden-1}, thutu=co]
\tieude{Người chơi đang dùng những phiên bản nào?}
\boicanh[Anh Duy · Kỹ sư QA]{Đội QA cần biết người chơi đang chạy những phiên bản nào và bản nào mới nhất để chọn bản kiểm thử hồi quy. Em liệt kê giúp anh, bản mới nhất ở trên cùng.}
\begin{yeucau}
  \item Xét các phiên chơi trong cửa sổ dữ liệu, từ ngày 1/7/2026 tới hết ngày 29/8/2026.
  \item Phiên bản \code{app_version} có dạng \code{'chính.phụ.vá'}. Tách ba phần thành số nguyên \code{so_chinh}, \code{so_phu}, \code{so_va}.
  \item Trả về \code{app_version}, \code{so_chinh}, \code{so_phu}, \code{so_va}, \code{so_phien}, \code{ngay_dau} (ngày sớm nhất có phiên của bản đó).
  \item Sắp theo \code{so_chinh}, rồi \code{so_phu}, rồi \code{so_va}, cả ba giảm dần.
\end{yeucau}
\begin{erd}
  \erdbang{sessions}{0,0}{session_id:PK, user_id:FK, event_date*, app_version*, platform}
  \erdbang{users}{1,0}{user_id:PK, install_date, app_version_install, platform}
  \erdnoi[kieu=n-1, tren={\code{sessions.user_id = users.user_id}}, duoi={nhiều phiên, một người chơi}]{sessions.user_id}{users.user_id}
\end{erd}
\begin{nghiepvu}{Xác định phiên bản cần kiểm thử hồi quy trước khi phát hành bản kế tiếp.}
  \khainiem{phien-ban-luc-choi}
  \luuy{So phiên bản theo chuỗi chỉ đúng khi mọi phần có cùng số chữ số; \code{'2.10.0'} đứng trước \code{'2.9.0'} nếu so theo chữ. Tách thành số rồi so từng phần mới đúng.}
\end{nghiepvu}
\goiy{Phần giữa nằm giữa dấu chấm thứ nhất và thứ hai, nên cần cắt hai lần.}
\goiy{\code{CAST(… AS UNSIGNED)} đổi chuỗi thành số nguyên; \code{SUBSTRING_INDEX(app_version, '.', -1)} lấy phần sau dấu chấm cuối.}
\ghichuthietke{Dễ: 2 bảng có điều kiện nối; câu chỉ cần \code{sessions}. Bài sai lấy \code{so_phu} bằng một lần cắt (\code{'2.1'} thành 2) rớt ở public; bài sai bỏ cửa sổ dữ liệu rớt ở hidden-1 (phiên 2026-09-01 bản 2.1.1). Dữ liệu chỉ có 2.0.2, 2.1.0, 2.1.1 nên \code{MAX} theo chuỗi vẫn đúng ở đây; đề yêu cầu tách số để học viên không dựa vào may mắn đó.}
\end{cauhoi}

% ================================================================ S6 · write · Dễ
\begin{cauhoi}[ma=jelly-s6-tuan-iso, dulieu=jelly, domain=customer,
  hinhthuc=write, cap=easy, dang=B1, kynang={YEARWEEK, WEEKDAY, DATE_SUB}, bay={tuan-bat-dau-chu-nhat, tuan-khong-tron},
  tacgia={Agent gói 2}, trangthai=review, bokiemtra={public, hidden-1}, thutu=co]
\tieude{Số người cài game theo tuần}
\boicanh[Chị Mai · Quản lý sản phẩm]{Chị cần biểu đồ số người chơi mới theo tuần cho buổi họp sản phẩm, tuần tính từ thứ Hai tới Chủ nhật. Em đếm giúp chị từ đầu tháng 7 tới nay.}
\begin{yeucau}
  \item Xét người chơi cài từ ngày 1/7/2026 tới hết ngày 29/8/2026.
  \item Tuần bắt đầu vào thứ Hai. \code{tuan} là số nguyên ghép năm và số thứ tự tuần theo chuẩn ISO, như \code{YEARWEEK} trả về (ví dụ số 202628 là tuần 28 năm 2026; không trả về chuỗi \code{'202628'}). \code{ngay_dau_tuan} là ngày thứ Hai của tuần, kể cả khi ngày đó nằm trước 1/7/2026.
  \item Trả về \code{tuan}, \code{ngay_dau_tuan}, \code{so_ngay_co_du_lieu} (số ngày trong tuần có người cài), \code{so_nguoi_cai}; sắp \code{tuan} tăng dần.
\end{yeucau}
\begin{erd}
  \erdbang{users}{0,0}{user_id:PK, install_date*, platform, country, channel}
  \erdbang{player_daily}{1,0}{user_id:FK, activity_date, install_date, cohort_day}
  \erdnoi[kieu=1-n, tren={\code{users.user_id = player_daily.user_id}}, duoi={một người chơi, nhiều ngày hoạt động}]{users.user_id}{player_daily.user_id}
\end{erd}
\begin{nghiepvu}{Theo dõi xu hướng người chơi mới theo tuần để đánh giá các đợt quảng cáo.}
  \khainiem{nhom-cai-dat}
  \luuy{Kỳ dữ liệu bắt đầu giữa tuần và kết thúc trước Chủ nhật, nên tuần đầu và tuần cuối không đủ 7 ngày. So các tuần với nhau cần nhìn cả số ngày có dữ liệu.}
\end{nghiepvu}
\goiy{\code{YEARWEEK} có tham số chế độ; chế độ mặc định bắt đầu tuần vào Chủ nhật. Chế độ 3 theo chuẩn ISO, tuần bắt đầu thứ Hai.}
\goiy{Ngày thứ Hai của tuần bằng ngày cài trừ đi \code{WEEKDAY(install_date)} ngày (\code{DATE_SUB}).}
\ghichuthietke{Dễ: 2 bảng có điều kiện nối; câu chỉ cần \code{users}. Hai bài sai rớt ở public: \code{YEARWEEK} chế độ mặc định (tuần từ Chủ nhật); lấy ngày cài sớm nhất làm ngày đầu tuần (tuần đầu ra 2026-07-01 thay vì thứ Hai 2026-06-29).}
\end{cauhoi}

% ================================================================ S7 · debug · Cơ bản
\begin{cauhoi}[ma=jelly-s7-thang-sai-thu-tu, dulieu=jelly, domain=finance,
  hinhthuc=debug, cap=basic, dang=B1, kynang={DATE_FORMAT, GROUP BY}, bay={sap-ten-thang},
  tacgia={Agent gói 2}, trangthai=review, bokiemtra={public, hidden-1}, thutu=co]
\tieude{Doanh thu IAP theo tháng sắp ngược thứ tự}
\boicanh[Bạn Vy · Thực tập sinh phân tích]{Em làm báo cáo doanh thu IAP (mua hàng trong ứng dụng) theo tháng mà tháng 8 lại đứng trước tháng 7. Em sửa mãi chưa được, anh chị chỉ giúp em với.}
\begin{yeucau}
  \item Truy vấn dưới đây tính doanh thu IAP theo tháng (giao dịch không hoàn tiền, từ 1/7/2026), nhưng tháng sắp sai thứ tự.
  \item Sửa để \code{thang} có dạng \code{'2026-07'} và các tháng sắp đúng thứ tự thời gian.
  \item Giữ nguyên điều kiện ngày, cột \code{thang}, \code{doanh_thu_iap_usd} (USD, làm tròn 2 chữ số).
\end{yeucau}
\begin{erd}
  \erdbang{purchases}{0,0}{transaction_id:PK, event_date*=2026-08-15, product_name='Season Pass', price_usd*=9.99, is_refund=0}
\end{erd}
\begin{nghiepvu}{Theo dõi doanh thu bán gói vật phẩm theo tháng.}
  \khainiem{doanh-thu-iap}
  \luuy{Tháng 8/2026 chỉ có dữ liệu tới hết ngày 29/8, nên doanh thu tháng 8 chưa trọn tháng.}
\end{nghiepvu}
\goiy{Tên tháng tiếng Anh sắp theo chữ cái: August đứng trước July. Mã định dạng nào cho tháng dạng hai chữ số?}
\ghichuthietke{Cơ bản: 1 bảng có giá trị mẫu; lỗi duy nhất ở mã định dạng \code{\%M}. Bài sai \code{'\%Y-\%M'} vẫn là tên tháng, rớt ở public. Bài sai bỏ điều kiện ngày (vì nghĩ dữ liệu chỉ có hai tháng) rớt ở hidden-1 (giao dịch 2026-09-01).}
\end{cauhoi}

% ================================================================ S8 · debug · Dễ
\begin{cauhoi}[ma=jelly-s8-gio-chua-doi, dulieu=jelly, domain=operations,
  hinhthuc=debug, cap=easy, dang=A1, kynang={DATETIME, khoảng nửa mở}, bay={between-tren-datetime},
  tacgia={Agent gói 2}, trangthai=review, bokiemtra={public, hidden-1}, thutu=co]
\tieude{Ca sáng 15/8 bị tính thừa một lượt sự cố}
\boicanh[Bạn Vy · Thực tập sinh phân tích]{Anh Duy bên QA cần các lượt chơi gặp sự cố sáng 15/8, tính tới trước 11 giờ là lúc đội QA giao ca. Danh sách em gửi có một lượt anh Duy bảo thuộc ca sau, anh chị xem lại giúp em.}
\begin{yeucau}
  \item Truy vấn dưới đây cần liệt kê lượt chơi có kết quả \code{'crash'} (mọi nền tảng) từ 0 giờ ngày 15/8/2026 tới trước 11 giờ 00 phút 00 giây cùng ngày. Lượt lúc đúng 11 giờ 00 phút 00 giây thuộc ca sau.
  \item Sửa điều kiện thời gian để danh sách chỉ còn các lượt của ca sáng; giữ nguyên cột \code{attempt_id}, \code{event_time}, \code{platform}, \code{level_id}, \code{app_version} và thứ tự sắp (\code{event_time} tăng dần, rồi \code{attempt_id}).
\end{yeucau}
\begin{erd}
  \erdbang{level_attempts}{0,0}{attempt_id:PK, level_id:FK, event_date, event_time*=2026-08-15 09:54:00, result*='crash', platform, app_version}
  \erdbang{level_metadata}{1,0}{level_id:PK, world_name, difficulty_tag}
  \erdnoi[kieu=n-1, tren={\code{level_attempts.level_id = level_metadata.level_id}}]{level_attempts.level_id}{level_metadata.level_id}
\end{erd}
\begin{nghiepvu}{Chia các lượt chơi gặp sự cố trong ngày phát hành bản 2.1.1 cho đúng ca QA để tái hiện lỗi.}
  \khainiem{phien-ban-luc-choi}
  \luuy{Cột thời điểm lưu cả ngày, giờ, phút và giây. Khoảng thời gian dạng ``từ mốc A tới trước mốc B'' không chứa mốc B.}
\end{nghiepvu}
\goiy{\code{BETWEEN} lấy cả hai đầu mút. Có lượt nào xảy ra đúng vào mốc cuối không?}
\goiy{Dùng khoảng nửa mở: \code{>=} mốc đầu và \code{<} mốc cuối.}
\ghichuthietke{Dễ: 2 bảng có điều kiện nối, giá trị mẫu có phần giây là manh mối. v2 (theo REVIEW-C): bỏ cách hỏi cả ngày vì dữ liệu không có lượt nào lúc đúng 0 giờ nên bài sai \code{BETWEEN} tới ngày hôm sau vẫn qua; đổi sang mốc 11 giờ vì có đúng một lượt sự cố (iOS) lúc 2026-08-15 11:00:00 ở cả public và hidden-1. Starter \code{BETWEEN … '2026-08-15 11:00:00'} thừa lượt đó. Bài sai \code{HOUR(event_time) <= 11} lấy cả giờ 11; bài sai chỉ có cận trên lấy cả các ngày trước; cả hai rớt ở public.}
\end{cauhoi}

% ================================================================ S9 · debug · Dễ
\begin{cauhoi}[ma=jelly-s9-ma-chien-dich, dulieu=jelly, domain=marketing,
  hinhthuc=debug, cap=easy, dang=A3, kynang={SUBSTRING, GROUP BY}, bay={ma-dai-ngan-khac-nhau},
  tacgia={Agent gói 2}, trangthai=review, bokiemtra={public, hidden-1}, thutu=co]
\tieude{Báo cáo người mới theo kênh có nhóm không tên}
\boicanh[Bạn Vy · Thực tập sinh phân tích]{Em đếm người chơi mới tháng 7 theo kênh bằng cách đọc mã chiến dịch, mà nhóm đông nhất lại không có tên kênh. Mai em phải nộp báo cáo rồi, em chưa biết sai ở bước nào.}
\begin{yeucau}
  \item Truy vấn dưới đây đếm người chơi cài trong tháng 7/2026 theo kênh thu hút, nhưng có một nhóm có \code{kenh} bằng NULL (không có tên kênh).
  \item Sửa để mọi người chơi có đúng tên kênh như cột \code{channel} (\code{'Organic'}, \code{'Search Ads'}, …).
  \item Giữ nguyên cột \code{kenh}, \code{so_nguoi_choi} và thứ tự sắp (\code{so_nguoi_choi} giảm dần, rồi \code{kenh}).
\end{yeucau}
\begin{erd}
  \erdbang{users}{0,0}{user_id:PK, install_date*, channel*='Organic', campaign_id*='ORG-000', platform}
  \erdbang{campaign_spend}{1,0}{date, campaign_id*='JTA-SEA-SEARCH-101', channel='Search Ads', spend_usd}
  \erdnoi[kieu=n-1, tren={\code{users.campaign_id = campaign_spend.campaign_id}}, duoi={nối logic, không có khoá ngoại}, canhbao]{users.campaign_id}{campaign_spend.campaign_id}
\end{erd}
\begin{nghiepvu}{Báo cáo số người chơi mới mỗi kênh mang về trong tháng.}
  \khainiem{kenh-thu-hut}
  \luuy{Mã chiến dịch của kênh trả phí và của kênh tự nhiên có độ dài, cấu trúc khác nhau. Tên kênh đã có sẵn ở một cột riêng.}
\end{nghiepvu}
\goiy{So giá trị mẫu \code{'ORG-000'} với \code{'JTA-SEA-SEARCH-101'}: ký tự thứ 5 tới thứ 7 của mỗi mã là gì?}
\goiy{Bảng người chơi đã có cột \code{channel}; gom theo cột này.}
\ghichuthietke{Dễ: 2 bảng, đường nối tô cảnh báo chỉ vào \code{campaign_id}. Starter cắt \code{SUBSTRING(campaign_id, 5, 3)} nên \code{'ORG-000'} ra \code{'000'}, kênh Organic thành NULL. Bài sai \code{SUBSTRING_INDEX(campaign_id, '-', 1)} ra \code{'JTA'} cho mọi kênh trả phí; bài sai thiếu cận trên lẫn tháng 8; cả hai rớt ở public.}
\end{cauhoi}


% ################################################################ MODULE 4 · AGGREGATION

% ================================================================ A1 · write · Dễ
\begin{cauhoi}[ma=jelly-a1-dau-thang-7, dulieu=jelly, domain=customer,
  hinhthuc=write, cap=easy, dang=A2, kynang={COUNT, COUNT DISTINCT}, bay={mau-chia-31, quen-can-tren},
  tacgia={Agent gói 2}, trangthai=review, bokiemtra={public, hidden-1}, thutu=co]
\tieude{DAU bình quân tháng 7 là bao nhiêu?}
\boicanh[Chị Mai · Quản lý sản phẩm]{Chị cần con số người chơi hoạt động mỗi ngày của tháng 7 để đặt mục tiêu cho quý sau. Em tính giúp chị bình quân một ngày có bao nhiêu người chơi.}
\begin{yeucau}
  \item Xét các ngày hoạt động từ 1/7/2026 tới hết 31/7/2026.
  \item Tính DAU bình quân của tháng bằng tổng DAU các ngày chia số ngày có dữ liệu.
  \item Trả về một dòng: \code{so_ngay}, \code{dau_tb} (người chơi, làm tròn 1 chữ số).
\end{yeucau}
\begin{erd}
  \erdbang{player_daily}{0,0}{user_id:FK, activity_date*, sessions, iap_revenue}
  \erdbang{users}{1,0}{user_id:PK, install_date, platform, country}
  \erdnoi[kieu=n-1, tren={\code{player_daily.user_id = users.user_id}}, duoi={nhiều ngày hoạt động, một người chơi}]{player_daily.user_id}{users.user_id}
\end{erd}
\begin{nghiepvu}{Đặt mục tiêu người chơi hoạt động hằng ngày cho quý sau.}
  \khainiem{dau}
  \luuy{Mỗi dòng của bảng hoạt động là một người chơi trong một ngày có mở game, nên số dòng của một ngày chính là DAU của ngày đó. Số người chơi khác nhau trong cả tháng là MAU (số người chơi hoạt động trong tháng), khác với tổng DAU.}
\end{nghiepvu}
\goiy{Tổng DAU của cả tháng bằng số dòng của tháng. Một người chơi hoạt động 20 ngày trong tháng được tính 20 lần.}
\goiy{\code{COUNT(*)} chia \code{COUNT(DISTINCT activity_date)}.}
\ghichuthietke{Dễ: 2 bảng có điều kiện nối; câu chỉ cần \code{player_daily}. Hai bài sai rớt ở public: đếm người chơi khác nhau cả tháng chia 31 (MAU chia số ngày); thiếu cận trên lẫn tháng 8. Học viên dùng \code{AVG(daily_kpi.dau)} cũng ra đúng.}
\end{cauhoi}

% ================================================================ A2 · write · Dễ
\begin{cauhoi}[ma=jelly-a2-doanh-thu-goi, dulieu=jelly, domain=sales,
  hinhthuc=write, cap=easy, dang=A3, kynang={GROUP BY, SUM, AVG}, bay={quen-cua-so, dem-nguoi-thay-giao-dich},
  tacgia={Agent gói 2}, trangthai=review, bokiemtra={public, hidden-1}, thutu=co]
\tieude{Doanh thu từng gói vật phẩm}
\boicanh[Chị Hương · Quản lý doanh thu]{Chị chuẩn bị đổi giá vài gói vật phẩm trong cửa hàng. Em tổng hợp giúp chị từ đầu tháng 7 tới nay mỗi gói bán được bao nhiêu lượt và mang về bao nhiêu tiền.}
\begin{yeucau}
  \item Xét các giao dịch không hoàn tiền trong cửa sổ dữ liệu, từ 1/7/2026 tới hết 29/8/2026.
  \item Tính cho từng gói vật phẩm (\code{product_name}) số giao dịch, doanh thu IAP và giá bình quân một giao dịch.
  \item Trả về \code{product_name}, \code{so_giao_dich}, \code{doanh_thu_usd} (USD, làm tròn 2 chữ số), \code{gia_tb_usd} (USD, làm tròn 2 chữ số).
  \item Sắp \code{doanh_thu_usd} giảm dần; bằng nhau thì \code{product_name} theo thứ tự chữ cái.
\end{yeucau}
\begin{erd}
  \erdbang{purchases}{0,0}{transaction_id:PK, user_id:FK, event_date*, product_name*, price_usd*, is_refund*}
  \erdbang{users}{1,0}{user_id:PK, install_date, platform, channel}
  \erdnoi[kieu=n-1, tren={\code{purchases.user_id = users.user_id}}, duoi={nhiều giao dịch, một người chơi}]{purchases.user_id}{users.user_id}
\end{erd}
\begin{nghiepvu}{Chọn gói vật phẩm cần đổi giá hoặc đổi vị trí trong cửa hàng.}
  \khainiem{doanh-thu-iap}
  \luuy{Doanh thu IAP chỉ tính giao dịch không hoàn tiền. Kỳ này chưa có giao dịch hoàn tiền, nhưng điều kiện vẫn cần có để báo cáo đúng khi phát sinh hoàn tiền.}
\end{nghiepvu}
\goiy{Mỗi dòng là một giao dịch; một người chơi có thể mua cùng một gói nhiều lần.}
\goiy{Lọc \code{is_refund = 0} và khoảng ngày có cả cận trên, rồi gom theo \code{product_name}.}
\ghichuthietke{Dễ: 2 bảng có điều kiện nối; câu chỉ cần \code{purchases}. Bài sai bỏ cửa sổ dữ liệu rớt ở hidden-1 (giao dịch Booster Bundle 2026-09-01); bài sai đếm người chơi khác nhau thay số giao dịch rớt ở public. Mỗi gói có giá cố định nên giá bình quân bằng giá niêm yết; cột này để học viên tự thấy điều đó.}
\end{cauhoi}

% ================================================================ A3 · write · Dễ
\begin{cauhoi}[ma=jelly-a3-chi-theo-kenh, dulieu=jelly, domain=marketing,
  hinhthuc=write, cap=easy, dang=A3, kynang={GROUP BY, SUM, COUNT DISTINCT}, bay={quen-can-tren, avg-theo-dong},
  tacgia={Agent gói 2}, trangthai=review, bokiemtra={public, hidden-1}, thutu=co]
\tieude{Chi phí quảng cáo tháng 8 của từng kênh}
\boicanh[Anh Tuấn · Quản lý thu hút người chơi]{Anh đang làm báo cáo chi phí quảng cáo tháng 8 gửi giám đốc studio. Em tính giúp anh tháng này mình đã chi cho mỗi kênh tổng cộng bao nhiêu, bình quân một ngày bao nhiêu.}
\begin{yeucau}
  \item Xét chi phí từ ngày 1/8/2026 tới hết tháng 8/2026 (dữ liệu tháng 8 có tới hết ngày 29/8).
  \item Tính cho từng kênh tổng chi phí và chi phí bình quân một ngày (tổng chi chia số ngày có chi phí).
  \item Trả về \code{channel}, \code{tong_chi_usd}, \code{chi_tb_ngay_usd} (USD, làm tròn 2 chữ số).
  \item Sắp \code{tong_chi_usd} giảm dần.
\end{yeucau}
\begin{erd}
  \erdbang{campaign_spend}{0,0}{date*, campaign_id, channel*, country, target_platform, spend_usd*}
  \erdbang{users}{1,0}{user_id:PK, install_date, channel, campaign_id}
  \erdnoi[kieu=1-n, tren={\code{campaign_spend.campaign_id = users.campaign_id}}, duoi={nối logic, không có khoá ngoại}]{campaign_spend.campaign_id}{users.campaign_id}
\end{erd}
\begin{nghiepvu}{Báo cáo chi phí quảng cáo theo kênh và chuẩn bị chia ngân sách tháng sau.}
  \khainiem{kenh-thu-hut}
  \luuy{Một ngày của một kênh có nhiều dòng chi phí, chia theo quốc gia và nền tảng nhắm tới. Bình quân một ngày phải chia cho số ngày, không chia cho số dòng.}
\end{nghiepvu}
\goiy{\code{AVG(spend_usd)} là chi bình quân của một dòng. Một ngày có mấy dòng cho mỗi kênh?}
\goiy{\code{SUM(spend_usd) / COUNT(DISTINCT `date`)}; lọc tháng có cả cận dưới và cận trên.}
\ghichuthietke{Dễ: 2 bảng nối logic; câu chỉ cần \code{campaign_spend}. Bài sai thiếu cận trên rớt ở hidden-1 (dòng Social Ads 2026-09-01); bài sai \code{AVG(spend_usd)} rớt ở public. Chọn tháng 8 để bẫy cận trên có tác dụng; đề ghi rõ tháng chưa trọn.}
\end{cauhoi}

% ================================================================ A4 · write · Trung bình
\begin{cauhoi}[ma=jelly-a4-ty-le-thang-man, dulieu=jelly, domain=operations,
  hinhthuc=write, cap=easy, dang=A2, kynang={SUM(CASE), GROUP BY, tỷ lệ}, bay={bo-crash-khoi-mau-so, quen-cua-so},
  tacgia={Agent gói 2}, trangthai=review, bokiemtra={public, hidden-1}, thutu=co, saiso=0.01]
\tieude{Tỷ lệ thắng của từng màn}
\boicanh[Anh Khoa · Thiết kế màn chơi]{Anh cần xếp lại độ khó của 60 màn cho mùa sau. Em tính giúp anh tỷ lệ thắng của từng màn từ đầu tháng 7 tới nay, màn khó nhất đứng đầu.}
\begin{yeucau}
  \item Xét các lượt chơi từ 1/7/2026 tới hết 29/8/2026.
  \item Tính tỷ lệ thắng của từng màn.
  \item Trả về \code{level_id}, \code{world_name}, \code{so_luot}, \code{ty_le_thang_pct} (phần trăm, làm tròn 2 chữ số).
  \item Sắp \code{ty_le_thang_pct} tăng dần; bằng nhau thì \code{level_id} nhỏ hơn đứng trước.
\end{yeucau}
\begin{erd}
  \erdbang{level_attempts}{0,0}{attempt_id:PK, level_id:FK, world_name*, event_date*, result*}
  \erdbang{level_metadata}{1,0}{level_id:PK, world_name, difficulty_tag}
  \erdnoi[kieu=n-1, tren={\code{level_attempts.level_id = level_metadata.level_id}}, duoi={nhiều lượt chơi, một màn}]{level_attempts.level_id}{level_metadata.level_id}
\end{erd}
\begin{nghiepvu}{Xếp lại độ khó các màn và chọn màn cần chỉnh trước mùa sau.}
  \khainiem{ty-le-thang-man}
  \luuy{Lượt gặp sự cố vẫn là một lượt người chơi đã vào màn. Bỏ các lượt này khỏi mẫu số làm tỷ lệ thắng cao hơn thực tế ở màn hay gặp sự cố.}
\end{nghiepvu}
\goiy{Mẫu số là mọi lượt chơi của màn, kể cả lượt thua, bỏ cuộc và gặp sự cố.}
\goiy{\code{SUM(CASE WHEN result = 'complete' THEN 1 ELSE 0 END)} chia \code{COUNT(*)}, nhân 100.}
\ghichuthietke{Dễ (v2 hạ từ Trung bình theo REVIEW-B: 1 bảng, 1 lần gom nhóm): 2 bảng có điều kiện nối; câu chỉ cần \code{level_attempts} (có sẵn \code{world_name}). Bài sai bỏ lượt \code{crash} khỏi mẫu số rớt ở public; bài sai bỏ cửa sổ dữ liệu rớt ở hidden-1 (lượt màn 58 ngày 2026-09-01).}
\end{cauhoi}

% ================================================================ A5 · write · Trung bình
\begin{cauhoi}[ma=jelly-a5-arppu, dulieu=jelly, domain=finance,
  hinhthuc=write, cap=medium, dang=A2, kynang={COUNT DISTINCT, SUM, GROUP BY}, bay={mau-so-giao-dich, quen-can-tren},
  tacgia={Agent gói 2}, trangthai=review, bokiemtra={public, hidden-1}, thutu=co, saiso=0.01]
\tieude{ARPPU tháng 7 là bao nhiêu?}
\boicanh[Chị Hương · Quản lý doanh thu]{Chị muốn biết mỗi người chịu trả tiền trên iOS và Android mang về bao nhiêu trong tháng 7, để định giá gói mới cho từng nền tảng.}
\begin{yeucau}
  \item Xét các giao dịch không hoàn tiền từ 1/7/2026 tới hết 31/7/2026.
  \item Tính ARPPU (doanh thu IAP bình quân trên một người chơi trả tiền) cho từng nền tảng của giao dịch (\code{purchases.platform}).
  \item Trả về \code{platform}, \code{doanh_thu_iap_usd} (USD, làm tròn 2 chữ số), \code{so_nguoi_tra_tien}, \code{arppu_usd} (USD, làm tròn 2 chữ số).
  \item Sắp \code{platform} theo thứ tự chữ cái.
\end{yeucau}
\begin{erd}
  \erdbang{purchases}{1,0}{transaction_id:PK, user_id:FK, session_id:FK, event_date, price_usd, platform}
  \erdbang{users}{0,0}{user_id:PK, platform}
  \erdbang{sessions}{2,0}{session_id:PK, user_id:FK}
  \erdnoi[kieu=1-n]{users.user_id}{purchases.user_id}
  \erdnoi[kieu=n-1]{purchases.session_id}{sessions.session_id}
\end{erd}
\begin{nghiepvu}{Định giá gói vật phẩm mới cho từng nền tảng.}
  \khainiem{arppu}
  \khainiem{nguoi-choi-tra-tien}
  \luuy{Một người chơi trả tiền có thể có nhiều giao dịch trong tháng. Chia doanh thu cho số giao dịch chỉ ra giá bình quân một giao dịch.}
\end{nghiepvu}
\goiy{Mẫu số của ARPPU đếm người, không đếm giao dịch.}
\goiy{\code{SUM(price_usd) / COUNT(DISTINCT user_id)}.}
\ghichuthietke{Trung bình: 3 bảng có ký hiệu, câu chỉ cần \code{purchases}. Bài sai chia cho \code{COUNT(user_id)} rớt ở public; bài sai thiếu cận trên lẫn tháng 8, rớt ở public.}
\end{cauhoi}

% ================================================================ A6 · write · Trung bình
\begin{cauhoi}[ma=jelly-a6-quoc-gia-nhieu-crash, dulieu=jelly, domain=after_sales,
  hinhthuc=write, cap=easy, dang=A3, kynang={SUM, tỷ lệ của tổng, GROUP BY}, bay={trung-binh-ty-le, quen-cua-so},
  tacgia={Agent gói 2}, trangthai=review, bokiemtra={public, hidden-1}, thutu=co, saiso=0.01]
\tieude{Thị trường nào có tỷ lệ sự cố cao nhất?}
\boicanh[Anh Duy · Kỹ sư QA]{Đội QA chỉ đủ người kiểm thử thêm trên thiết bị của hai thị trường. Hai tháng qua thị trường nào có tỷ lệ sự cố cao nhất, em gửi anh trước chiều mai nhé.}
\begin{yeucau}
  \item Xét các ngày hoạt động từ 1/7/2026 tới hết 29/8/2026.
  \item Tính cho từng thị trường (\code{country}) số sự cố, số phiên và tỷ lệ sự cố trên 1.000 phiên (công thức ở tab Nghiệp vụ).
  \item Trả về \code{country}, \code{so_su_co}, \code{so_phien}, \code{su_co_tren_1000_phien} (làm tròn 2 chữ số).
  \item Sắp \code{su_co_tren_1000_phien} giảm dần; bằng nhau thì \code{country} theo thứ tự chữ cái.
\end{yeucau}
\begin{erd}
  \erdbang{player_daily}{0,0}{user_id:FK, activity_date*, country*, sessions*, crash_count*}
  \erdbang{users}{1,0}{user_id:PK, country, platform}
  \erdnoi[kieu=n-1, tren={\code{player_daily.user_id = users.user_id}}, duoi={nhiều ngày hoạt động, một người chơi}]{player_daily.user_id}{users.user_id}
\end{erd}
\begin{nghiepvu}{Ưu tiên thị trường cho đợt kiểm thử thêm trên thiết bị thật.}
  \khainiem{ty-le-su-co}
  \luuy{Tỷ lệ của một thị trường tính từ tổng sự cố chia tổng phiên. Nếu lấy trung bình tỷ lệ của từng dòng (một người chơi trong một ngày), ngày chỉ có một phiên được tính nặng ngang ngày có nhiều phiên.}
\end{nghiepvu}
\goiy{Cộng sự cố và cộng phiên của cả thị trường trước, rồi mới chia.}
\goiy{\code{1000 * SUM(crash_count) / SUM(sessions)}.}
\ghichuthietke{Dễ (v2 hạ từ Trung bình theo REVIEW-B): 2 bảng có điều kiện nối; câu chỉ cần \code{player_daily}. Số lượt \code{crash} ở \code{level_attempts} là nguồn đếm khác, không khớp \code{crash_count}. Bài sai trung bình tỷ lệ từng dòng rớt ở public; bài sai bỏ cửa sổ dữ liệu rớt ở hidden-1 (dòng GB ngày 2026-09-01).}
\end{cauhoi}

% ================================================================ A7 · debug · Dễ
\begin{cauhoi}[ma=jelly-a7-dem-nguoi-tra-tien, dulieu=jelly, domain=sales,
  hinhthuc=debug, cap=easy, dang=A2, kynang={COUNT DISTINCT}, bay={dem-giao-dich-thay-nguoi},
  tacgia={Agent gói 2}, trangthai=review, bokiemtra={public, hidden-1}, thutu=co]
\tieude{Số người trả tiền tháng 8 cao bất thường}
\boicanh[Chị Hương · Quản lý doanh thu]{Báo cáo số người trả tiền tháng 8 theo nền tảng ra cao hơn số bên cửa hàng ứng dụng gửi. Em kiểm tra giúp chị.}
\begin{yeucau}
  \item Truy vấn dưới đây cần đếm số người chơi trả tiền, không phải số giao dịch, của từng nền tảng trong tháng 8/2026 (dữ liệu tới hết ngày 29/8).
  \item Sửa truy vấn để mỗi người chơi chỉ được đếm một lần; giữ nguyên cột \code{platform}, \code{so_nguoi_tra_tien}, điều kiện lọc và thứ tự sắp.
\end{yeucau}
\begin{erd}
  \erdbang{users}{0,0}{user_id:PK, platform, install_date}
  \erdbang{purchases}{1,0}{transaction_id:PK, user_id:FK*, event_date*, platform*, price_usd}
  \erdnoi[kieu=1-n, tren={một người chơi, nhiều giao dịch}, canhbao]{users.user_id}{purchases.user_id}
\end{erd}
\begin{nghiepvu}{Đối chiếu số người trả tiền với báo cáo của cửa hàng ứng dụng.}
  \khainiem{nguoi-choi-tra-tien}
  \luuy{Một người chơi có thể mua nhiều lần trong tháng. Người trả tiền được xác định từ giao dịch, mỗi người chỉ tính một lần.}
\end{nghiepvu}
\goiy{Mỗi dòng của \code{purchases} là một giao dịch. Cần đếm gì để mỗi người chơi chỉ tính một lần?}
\ghichuthietke{Dễ: nhãn trên đường nối tô cảnh báo chỉ ra quan hệ một người chơi nhiều giao dịch. Bài sai đếm mã giao dịch khác nhau rớt ở public; bài sai bỏ cận trên rớt ở hidden-1 (người chơi Android mua ngày 2026-09-01, chưa mua trong tháng 8).}
\end{cauhoi}

% ================================================================ A8 · debug · Trung bình
\begin{cauhoi}[ma=jelly-a8-tb-ty-le, dulieu=jelly, domain=finance,
  hinhthuc=debug, cap=medium, dang=A2, kynang={SUM, tỷ lệ của tổng}, bay={trung-binh-ty-le},
  tacgia={Agent gói 2}, trangthai=review, bokiemtra={public, hidden-1}, thutu=co, saiso=0.00001]
\tieude{ARPDAU cả tháng 7 bị sai}
\boicanh[Chị Lan · Kế toán quản trị]{ARPDAU (doanh thu bình quân trên một người chơi hoạt động trong ngày) tháng 7 trong báo cáo không khớp số chị tự đối chiếu. Em tìm giúp chị chỗ sai.}
\begin{yeucau}
  \item Truy vấn dưới đây tính tổng doanh thu (doanh thu IAP từ mua hàng trong ứng dụng cộng doanh thu quảng cáo), tổng DAU và ARPDAU của cả tháng 7/2026 từ bảng chỉ số ngày, nhưng ARPDAU của tháng ra sai.
  \item Sửa cột \code{arpdau} cho đúng; giữ nguyên cột \code{tong_doanh_thu_usd}, \code{tong_dau}; \code{arpdau} tính bằng USD, làm tròn 4 chữ số.
\end{yeucau}
\begin{erd}
  \erdbang{daily_kpi}{0,0}{date:PK, dau, iap_revenue, ad_revenue, arpdau}
  \erdbang{player_daily}{1,0}{user_id:FK, activity_date, iap_revenue, ad_revenue}
  \erdbang{users}{2,0}{user_id:PK, install_date}
  \erdnoi[kieu=1-n]{daily_kpi.date}{player_daily.activity_date}
  \erdnoi[kieu=n-1]{player_daily.user_id}{users.user_id}
\end{erd}
\begin{nghiepvu}{Báo cáo ARPDAU tháng cho ban giám đốc studio.}
  \khainiem{dau}
  \khainiem{doanh-thu-quang-cao}
  \luuy{ARPDAU của một ngày bằng tổng doanh thu IAP và doanh thu quảng cáo của ngày đó chia DAU của ngày đó. Bảng chỉ số ngày đã có sẵn ARPDAU của từng ngày, làm tròn 4 chữ số.}
\end{nghiepvu}
\goiy{ARPDAU của nhiều ngày là tỷ lệ của tổng hay trung bình của tỷ lệ?}
\goiy{Tổng doanh thu IAP cộng quảng cáo chia tổng \code{dau}.}
\ghichuthietke{Trung bình: 3 bảng có ký hiệu, câu chỉ cần \code{daily_kpi}. v2 (theo REVIEW-B): bỏ khái niệm \code{arpdau} và lưu ý cũ khỏi tab vì hai chỗ này ghi thẳng ``không lấy trung bình cột arpdau''; tab chỉ còn định nghĩa ARPDAU một ngày. Starter lấy \code{AVG(arpdau)}, lệch khoảng 0,0035 USD nên \code{saiso} đặt 0,00001. Hai bài sai rớt ở public: quên doanh thu quảng cáo ở tử số; chia cho số ngày.}
\end{cauhoi}

% ================================================================ A9 · debug · Trung bình
\begin{cauhoi}[ma=jelly-a9-having-where, dulieu=jelly, domain=operations,
  hinhthuc=debug, cap=medium, dang=A3, kynang={WHERE, HAVING, tỷ lệ}, bay={loc-dong-truoc-khi-gom},
  tacgia={Agent gói 2}, trangthai=review, bokiemtra={public, hidden-1}, thutu=co, saiso=0.01]
\tieude{Lọc màn khó mà cả 60 màn đều lọt}
\boicanh[Anh Khoa · Thiết kế màn chơi]{Anh nhờ bên phân tích lọc các màn có tỷ lệ thua trên 25\%, nhưng danh sách ra đủ cả 60 màn, màn nào cũng thua 100\%. Em sửa giúp anh.}
\begin{yeucau}
  \item Truy vấn dưới đây cần liệt kê các màn có tỷ lệ thua do hết lượt đi trên 25\% trong toàn bộ dữ liệu. Tỷ lệ thua là số lượt \code{result} = \code{'fail_out_of_moves'} chia tổng số lượt chơi của màn, kể cả lượt thắng, bỏ cuộc và gặp sự cố.
  \item Sửa truy vấn để tỷ lệ thua tính trên mọi lượt chơi của màn; giữ nguyên cột \code{level_id}, \code{so_luot} (tổng số lượt chơi của màn), \code{ty_le_thua_pct} (phần trăm, làm tròn 2 chữ số), ngưỡng 25 và thứ tự sắp (\code{ty_le_thua_pct} giảm dần, rồi \code{level_id}).
\end{yeucau}
\begin{erd}
  \erdbang{level_attempts}{1,0}{attempt_id:PK, level_id:FK, session_id:FK, result, moves_left}
  \erdbang{level_metadata}{0,0}{level_id:PK, world_name, difficulty_tag}
  \erdbang{booster_events}{2,0}{booster_event_id:PK, attempt_id:FK}
  \erdnoi[kieu=1-n]{level_metadata.level_id}{level_attempts.level_id}
  \erdnoi[kieu=1-n]{level_attempts.attempt_id}{booster_events.attempt_id}
\end{erd}
\begin{nghiepvu}{Tìm màn cần giảm độ khó hoặc thêm lượt đi.}
  \luuy{Số lượt thua phụ thuộc số người chơi tới được màn, nên màn đầu game luôn có nhiều lượt thua. Muốn so độ khó giữa các màn cần tỷ lệ thua. Thua do hết lượt đi khác bỏ cuộc và gặp sự cố.}
\end{nghiepvu}
\goiy{Điều kiện ở \code{WHERE} bỏ dòng trước khi gom nhóm. Sau khi lọc như truy vấn hiện tại, mẫu số của tỷ lệ còn những lượt nào?}
\goiy{Bỏ điều kiện loại kết quả ở \code{WHERE}; đếm lượt thua có điều kiện trong tử số, giữ ngưỡng ở \code{HAVING}.}
\ghichuthietke{Trung bình: 3 bảng có ký hiệu, \code{booster_events} là bảng nhiễu. v2 (theo REVIEW-C): đổi từ số lượt thua tuyệt đối (dồn về màn đầu, đề v1 ghi sai ``gần hết các màn'' trong khi starter ra 37/60) sang tỷ lệ thua. Starter lọc \code{result} ở \code{WHERE} nên mẫu số chỉ còn lượt thua, màn nào cũng 100\% và ra đủ 60 màn. Bài sai bỏ lượt sự cố khỏi mẫu số, bài sai gộp cả bỏ cuộc vào tử số; cả hai rớt ở public. Màn 60 có rất ít lượt chơi; đề không đặt ngưỡng số lượt tối thiểu để giữ câu một lỗi.}
\end{cauhoi}


% ################################################################ CASE STUDY

% ================================================================ C3 · case · Cơ bản
\begin{cauhoi}[ma=jelly-c3-phien-ban-moi-dai-hon, dulieu=jelly, domain=operations,
  hinhthuc=case, cap=basic, diem=10, tacgia={Agent gói 2}, trangthai=review, bokiemtra={public, hidden-1}, thutu=co]
\tieude{Phiên chơi có dài hơn sau bản cập nhật?}
\boicanh[Chị Mai · Quản lý sản phẩm]{Đội sản phẩm nói từ bản 2.1.0 người chơi ở lại mỗi lần mở game lâu hơn. Em xem số liệu rồi cho chị kết luận trước buổi họp.}
\begin{erd}
  \erdbang{sessions}{0,0}{session_id:PK, event_date*=2026-08-15, duration_minutes*=12, app_version*='2.1.0', entry_screen}
\end{erd}
\begin{nghiepvu}{Quyết định có ghi vào báo cáo sản phẩm rằng bản cập nhật làm phiên chơi dài hơn hay không.}
  \khainiem{phien-choi}
  \khainiem{phien-ban-luc-choi}
  \luuy{Mỗi phiên bản bắt đầu có từ một ngày khác nhau, nên các bản được đo trên những khoảng ngày dài ngắn khác nhau.}
\end{nghiepvu}
\begin{buoc}[ma=b1, diem=5, hinhthuc=write, bang={sessions}]{Thời lượng phiên theo phiên bản}
  \begin{yeucau}
    \item Xét các phiên từ 1/7/2026 tới hết 29/8/2026, theo phiên bản lúc chơi.
    \item Trả về \code{app_version}, \code{so_phien}, \code{thoi_luong_tb} (thời lượng phiên bình quân, phút, làm tròn 2 chữ số), \code{ngay_dau}, \code{ngay_cuoi} (ngày sớm nhất và muộn nhất có phiên của bản đó).
    \item Sắp từ bản cũ tới bản mới.
  \end{yeucau}
  \goiy{Phiên bản lúc chơi nằm ngay trong bảng phiên; lọc ngày có cả cận trên.}
\end{buoc}
\begin{buoc}[ma=b2, diem=5, hinhthuc=ketluan, bang={sessions}]{Kết luận}
  \begin{luachon}
    \phuongan{Bản 2.1.1 có thời lượng phiên cao hơn bản 2.1.0, nên bản vá đã làm người chơi ở lại lâu hơn mỗi lần mở game.}
      {Chưa đúng. Chênh lệch giữa hai bản rất nhỏ, và hai bản được đo trên những khoảng ngày khác nhau.}
    \phuongan*{Thời lượng phiên bình quân của ba phiên bản gần như bằng nhau, nên chưa có căn cứ nói phiên chơi dài hơn sau bản cập nhật.}
      {Đúng. Bước 1 cho thấy ba bản chênh nhau chưa tới 0,1 phút, và mỗi bản chạy trong một khoảng ngày khác nhau.}
    \phuongan{Bản 2.0.2 có nhiều phiên nhất và phiên dài nhất, nên các bản cập nhật đã làm người chơi chơi ngắn đi.}
      {Chưa đúng. Bản 2.0.2 nhiều phiên vì có mặt từ đầu kỳ; mức chênh thời lượng không đủ lớn để kết luận.}
  \end{luachon}
\end{buoc}
\ghichuthietke{Case Cơ bản: 1 bảng có giá trị mẫu, 1 bước SQL, 1 bước kết luận. Phương án sai đọc chênh lệch rất nhỏ thành tác động của bản cập nhật, hoặc lấy số phiên tuyệt đối. Bài sai bước 1 bỏ cửa sổ dữ liệu rớt ở hidden-1 (phiên 2.1.1 ngày 2026-09-01 làm đổi \code{ngay_cuoi}); bài sai lấy phiên bản lúc cài rớt ở public.}
\end{cauhoi}

% ================================================================ C4 · case · Dễ
\begin{cauhoi}[ma=jelly-c4-kenh-dong-nguoi, dulieu=jelly, domain=marketing,
  hinhthuc=case, cap=easy, diem=20, tacgia={Agent gói 2}, trangthai=review, bokiemtra={public, hidden-1}, thutu=co]
\tieude{Kênh nào mang về nhiều người chơi trả tiền?}
\boicanh[Anh Tuấn · Quản lý thu hút người chơi]{Quý sau anh phải chia lại ngân sách giữa các kênh quảng cáo. Em xem người chơi cài tháng 7 của từng kênh có chịu trả tiền không rồi cho anh kết luận.}
\begin{erd}
  \erdbang{users}{0,0}{user_id:PK, install_date*, channel*, campaign_id}
  \erdbang{purchases}{1,0}{transaction_id:PK, user_id:FK, event_date*, price_usd, is_refund}
  \erdnoi[kieu=1-n, tren={\code{users.user_id = purchases.user_id}}, duoi={một người chơi, nhiều giao dịch}]{users.user_id}{purchases.user_id}
\end{erd}
\begin{nghiepvu}{Chia ngân sách quảng cáo quý sau giữa các kênh trả phí.}
  \khainiem{kenh-thu-hut}
  \khainiem{nguoi-choi-tra-tien}
  \luuy{Số người cài của các kênh chênh nhau nhiều. Kênh tự nhiên không có chi phí, nên không thể tăng hay giảm ngân sách cho kênh này.}
\end{nghiepvu}
\begin{buoc}[ma=b1, diem=8, hinhthuc=write, bang={users}]{Số người cài tháng 7 theo kênh}
  \begin{yeucau}
    \item Xét người chơi cài từ 1/7/2026 tới hết 31/7/2026.
    \item Trả về \code{channel}, \code{so_nguoi_cai}; sắp \code{so_nguoi_cai} giảm dần, bằng nhau thì \code{channel} theo thứ tự chữ cái.
  \end{yeucau}
  \goiy{Lọc ngày cài có cả cận dưới và cận trên.}
\end{buoc}
\begin{buoc}[ma=b2, diem=6, hinhthuc=write, bang={users, purchases}]{Tỷ lệ người chơi trả tiền theo kênh}
  \begin{yeucau}
    \item Cùng nhóm người chơi của bước 1. Người chơi trả tiền là người có ít nhất một giao dịch không hoàn tiền từ ngày cài tới hết 29/8/2026.
    \item Trả về \code{channel}, \code{so_nguoi_cai}, \code{so_nguoi_tra_tien}, \code{ty_le_tra_tien_pct} (phần trăm, làm tròn 2 chữ số).
    \item Sắp \code{ty_le_tra_tien_pct} giảm dần; bằng nhau thì \code{channel} theo thứ tự chữ cái.
  \end{yeucau}
  \goiy{Người chưa mua vẫn thuộc mẫu số; mỗi người trả tiền chỉ tính một lần dù mua nhiều lần.}
\end{buoc}
\begin{buoc}[ma=b3, diem=6, hinhthuc=ketluan, bang={users, purchases}]{Kết luận}
  \begin{luachon}
    \phuongan{Kênh tự nhiên (Organic) mang về nhiều người chơi trả tiền nhất, nên cần dồn thêm ngân sách cho kênh này, vì quy mô kênh đủ lớn để con số ổn định.}
      {Chưa đúng. Kênh tự nhiên đông nhất vì quy mô lớn, và kênh này không có chi phí để tăng ngân sách.}
    \phuongan*{Kênh tự nhiên (Organic) có nhiều người trả tiền nhất về số lượng; xét tỷ lệ thì Creator Campaign cao nhất, Video Network thấp nhất, nhưng Creator Campaign ít người cài nên mức chênh vài điểm phần trăm chưa chắc ổn định.}
      {Đúng. Bước 2 cho thấy kênh tự nhiên có nhiều người trả tiền nhất, còn xếp theo tỷ lệ thì Creator Campaign đứng đầu, Video Network đứng cuối; bước 1 cho thấy Creator Campaign có ít người cài nhất.}
    \phuongan{Search Ads và Social Ads mang về nhiều người cài nhất trong các kênh trả phí, nên cũng có tỷ lệ trả tiền cao nhất và chắc chắn đáng giữ ngân sách.}
      {Chưa đúng. Nhiều người cài không có nghĩa là tỷ lệ trả tiền cao; bước 2 cho thấy Creator Campaign ít người cài nhất nhưng tỷ lệ cao nhất.}
  \end{luachon}
\end{buoc}
\ghichuthietke{Case Dễ: ERD cố định 2 bảng có điều kiện nối; bước 1 chỉ dùng \code{users}, bước 2 tô sáng thêm \code{purchases}. Chọn người cài tháng 7 để mọi người chơi có ít nhất 29 ngày quan sát. Bài sai bước 1 thiếu cận trên rớt ở public; bài sai bước 2 đếm giao dịch thay người và sắp theo số tuyệt đối, đều rớt ở public. Kết luận chưa xét chi phí và doanh thu; phần đó để case c6.}
\end{cauhoi}

\end{document}
